% Options for packages loaded elsewhere
\PassOptionsToPackage{unicode}{hyperref}
\PassOptionsToPackage{hyphens}{url}
%
\documentclass[
]{article}
\usepackage{amsmath,amssymb}
\usepackage{iftex}
\ifPDFTeX
  \usepackage[T1]{fontenc}
  \usepackage[utf8]{inputenc}
  \usepackage{textcomp} % provide euro and other symbols
\else % if luatex or xetex
  \usepackage{unicode-math} % this also loads fontspec
  \defaultfontfeatures{Scale=MatchLowercase}
  \defaultfontfeatures[\rmfamily]{Ligatures=TeX,Scale=1}
\fi
\usepackage{lmodern}
\ifPDFTeX\else
  % xetex/luatex font selection
\fi
% Use upquote if available, for straight quotes in verbatim environments
\IfFileExists{upquote.sty}{\usepackage{upquote}}{}
\IfFileExists{microtype.sty}{% use microtype if available
  \usepackage[]{microtype}
  \UseMicrotypeSet[protrusion]{basicmath} % disable protrusion for tt fonts
}{}
\makeatletter
\@ifundefined{KOMAClassName}{% if non-KOMA class
  \IfFileExists{parskip.sty}{%
    \usepackage{parskip}
  }{% else
    \setlength{\parindent}{0pt}
    \setlength{\parskip}{6pt plus 2pt minus 1pt}}
}{% if KOMA class
  \KOMAoptions{parskip=half}}
\makeatother
\usepackage{xcolor}
\usepackage[margin=1in]{geometry}
\usepackage{longtable,booktabs,array}
\usepackage{calc} % for calculating minipage widths
% Correct order of tables after \paragraph or \subparagraph
\usepackage{etoolbox}
\makeatletter
\patchcmd\longtable{\par}{\if@noskipsec\mbox{}\fi\par}{}{}
\makeatother
% Allow footnotes in longtable head/foot
\IfFileExists{footnotehyper.sty}{\usepackage{footnotehyper}}{\usepackage{footnote}}
\makesavenoteenv{longtable}
\usepackage{graphicx}
\makeatletter
\newsavebox\pandoc@box
\newcommand*\pandocbounded[1]{% scales image to fit in text height/width
  \sbox\pandoc@box{#1}%
  \Gscale@div\@tempa{\textheight}{\dimexpr\ht\pandoc@box+\dp\pandoc@box\relax}%
  \Gscale@div\@tempb{\linewidth}{\wd\pandoc@box}%
  \ifdim\@tempb\p@<\@tempa\p@\let\@tempa\@tempb\fi% select the smaller of both
  \ifdim\@tempa\p@<\p@\scalebox{\@tempa}{\usebox\pandoc@box}%
  \else\usebox{\pandoc@box}%
  \fi%
}
% Set default figure placement to htbp
\def\fps@figure{htbp}
\makeatother
\setlength{\emergencystretch}{3em} % prevent overfull lines
\providecommand{\tightlist}{%
  \setlength{\itemsep}{0pt}\setlength{\parskip}{0pt}}
\setcounter{secnumdepth}{-\maxdimen} % remove section numbering
\usepackage{bookmark}
\IfFileExists{xurl.sty}{\usepackage{xurl}}{} % add URL line breaks if available
\urlstyle{same}
\hypersetup{
  pdftitle={hake\_parameter\_mapping},
  pdfauthor={Jon Brodziak},
  hidelinks,
  pdfcreator={LaTeX via pandoc}}

\title{hake\_parameter\_mapping}
\author{Jon Brodziak}
\date{2026-02-27}

\begin{document}
\maketitle

\section{\texorpdfstring{Current parameter mapping for \texttt{hake.R}
(driven by
\texttt{hake.inp})}{Current parameter mapping for hake.R (driven by hake.inp)}}\label{current-parameter-mapping-for-hake.r-driven-by-hake.inp}

This document maps \textbf{input parameters} read from \texttt{hake.inp}
into the \textbf{objects and roles} they play in \texttt{hake.R}, and
then maps the \textbf{estimated parameters} for models \textbf{(i)--(v)}
to the \textbf{exact optimization blocks} in the script.

\begin{quote}
Updated: 2026-02-27 (based on the current \texttt{hake.R} in this
project)
\end{quote}

\begin{center}\rule{0.5\linewidth}{0.5pt}\end{center}

\subsection{\texorpdfstring{1) How \texttt{hake.inp} is
read}{1) How hake.inp is read}}\label{how-hake.inp-is-read}

\texttt{hake.R} reads \texttt{hake.inp} as \textbf{key=value} pairs via
\texttt{read.table(...,\ sep="=")} and stores them in \texttt{kv} (a
named character vector). Typed parameters are then created with
\texttt{as.integer()} / \texttt{as.numeric()} conversions.

\begin{itemize}
\tightlist
\item
  Input parsing: around lines \textbf{236--243}
\item
  Typed parameter assignment: around lines \textbf{244--290} (see table
  below)
\end{itemize}

\begin{center}\rule{0.5\linewidth}{0.5pt}\end{center}

\subsection{\texorpdfstring{2) Input parameters from
\texttt{hake.inp}}{2) Input parameters from hake.inp}}\label{input-parameters-from-hake.inp}

\subsubsection{Summary table (keys → R objects →
meaning)}\label{summary-table-keys-r-objects-meaning}

\begin{longtable}[]{@{}
  >{\raggedright\arraybackslash}p{(\linewidth - 8\tabcolsep) * \real{0.1765}}
  >{\raggedleft\arraybackslash}p{(\linewidth - 8\tabcolsep) * \real{0.2353}}
  >{\raggedleft\arraybackslash}p{(\linewidth - 8\tabcolsep) * \real{0.2353}}
  >{\raggedright\arraybackslash}p{(\linewidth - 8\tabcolsep) * \real{0.1765}}
  >{\raggedright\arraybackslash}p{(\linewidth - 8\tabcolsep) * \real{0.1765}}@{}}
\toprule\noalign{}
\begin{minipage}[b]{\linewidth}\raggedright
\texttt{hake.inp} key
\end{minipage} & \begin{minipage}[b]{\linewidth}\raggedleft
R object in \texttt{hake.R}
\end{minipage} & \begin{minipage}[b]{\linewidth}\raggedleft
Type
\end{minipage} & \begin{minipage}[b]{\linewidth}\raggedright
Used for
\end{minipage} & \begin{minipage}[b]{\linewidth}\raggedright
Where used in \texttt{hake.R} (typical)
\end{minipage} \\
\midrule\noalign{}
\endhead
\bottomrule\noalign{}
\endlastfoot
\texttt{K} & \texttt{K} & integer & Number of composition categories
(simplex dimension) & mesh construction; likelihood dimensions (e.g.,
\texttt{mk\_simplex(K,h)}, \texttt{ncol=K}) \\
\texttt{G} & \texttt{G} & integer & Number of independent samples /
groups & simulation loop over groups; groupwise likelihood sums \\
\texttt{h} & \texttt{h} & numeric & Grid step for
\texttt{mk\_simplex(K,h)} (mesh resolution) & controls number of simplex
points \texttt{nmesh} \\
\texttt{theta\_true} & \texttt{theta\_true} & numeric & ``True'' linear
Dirichlet--multinomial scaling used in simulation; also used as fixed θ
in model (iii) & simulation:
\texttt{a0g\ \textless{}-\ (theta\_true*od\_mult)*N\_row{[}g{]}}; model
(iii): \texttt{theta\_fixed\ \textless{}-\ theta\_true} \\
\texttt{od\_mult} & \texttt{od\_mult} & numeric & Overdispersion
multiplier applied to α₀,g (smaller → more overdispersion) & simulation
only: multiplies α₀,g before drawing Dirichlet \\
\texttt{sigma} & \texttt{sigma} & numeric & Between-sample variation for
perturbed group proportions \texttt{p\_g} & simulation only:
\texttt{z\ \textless{}-\ log(p\_true)\ +\ rnorm(K,0,sigma)} \\
\texttt{nsims} & \texttt{nsims} & integer & Replicates per simplex point
(Monte Carlo) & outer simulation loop / repeated RMSE estimation \\
\texttt{random.seed} & \texttt{random.seed} & integer & RNG seed &
\texttt{set.seed(random.seed)} \\
\texttt{dist\_code} & \texttt{dist\_code} & integer & Sample-size
generator selection (1--4) & drives \texttt{draw\_nsamp()} \\
\texttt{mean\_nsamp} & \texttt{mean\_nsamp} & numeric & Target mean
sample size across distributions & used by all \texttt{dist\_code}
cases \\
\texttt{ln\_sd} & \texttt{ln\_sd} & numeric & Lognormal sdlog (only when
\texttt{dist\_code==2}) & \texttt{draw\_nsamp()} lognormal case \\
\texttt{nb\_size} & \texttt{nb\_size} & numeric & NegBin ``size'' (only
when \texttt{dist\_code==3}) & \texttt{draw\_nsamp()} negbin case \\
\texttt{Nmin} & \texttt{Nmin} & integer & Baseline lower bound for
log-uniform sample sizes (only when \texttt{dist\_code==4}) &
\texttt{draw\_nsamp()} log-uniform case (scaled to hit
\texttt{mean\_nsamp}) \\
\texttt{Nmax} & \texttt{Nmax} & integer & Baseline upper bound for
log-uniform sample sizes (only when \texttt{dist\_code==4}) &
\texttt{draw\_nsamp()} log-uniform case (scaled to hit
\texttt{mean\_nsamp}) \\
\texttt{p\_ubound} & \texttt{p\_ubound} & numeric vector length
\texttt{K} & Component-wise \textbf{upper} bounds for filtering simplex
points & mesh filter: keep rows where
\texttt{p\_true{[}k{]}\ \textless{}=\ p\_ubound{[}k{]}} \\
\texttt{p\_lbound} & \texttt{p\_lbound} & numeric vector length
\texttt{K} & Component-wise \textbf{lower} bounds for filtering simplex
points & mesh filter: keep rows where
\texttt{p\_true{[}k{]}\ \textgreater{}=\ p\_lbound{[}k{]}} \\
\end{longtable}

\subsubsection{Required vs conditional
keys}\label{required-vs-conditional-keys}

\begin{itemize}
\item
  \textbf{Always required} (must be present and numeric):\\
  \texttt{K,\ G,\ h,\ theta\_true,\ dist\_code,\ mean\_nsamp,\ nsims,\ random.seed,\ sigma,\ od\_mult}
\item
  \textbf{Conditionally required}:

  \begin{itemize}
  \tightlist
  \item
    If \texttt{dist\_code==2} (lognormal): \texttt{ln\_sd}
  \item
    If \texttt{dist\_code==3} (negative binomial): \texttt{nb\_size}
  \item
    If \texttt{dist\_code==4} (log-uniform): \texttt{Nmin,\ Nmax}
  \end{itemize}
\item
  \textbf{Optional}:

  \begin{itemize}
  \tightlist
  \item
    \texttt{p\_ubound} defaults to \texttt{rep(1,\ K)} if blank/missing
  \item
    \texttt{p\_lbound} defaults to \texttt{rep(0,\ K)} if blank/missing
  \end{itemize}
\end{itemize}

\begin{center}\rule{0.5\linewidth}{0.5pt}\end{center}

\subsection{3) Derived objects: where inputs turn into data used for
estimation}\label{derived-objects-where-inputs-turn-into-data-used-for-estimation}

\subsubsection{3.1 Mesh of ``true''
compositions}\label{mesh-of-true-compositions}

\begin{itemize}
\tightlist
\item
  \texttt{mesh\ \textless{}-\ mk\_simplex(K,\ h)} creates a grid of
  simplex points \texttt{p\_true} (one per row).
\item
  After optional filtering via \texttt{p\_lbound}/\texttt{p\_ubound},
  the script loops over rows of \texttt{mesh}.
\end{itemize}

\subsubsection{\texorpdfstring{3.2 Sample sizes per group
(\texttt{N\_vec})}{3.2 Sample sizes per group (N\_vec)}}\label{sample-sizes-per-group-n_vec}

\begin{itemize}
\tightlist
\item
  \texttt{N\_vec} is an \texttt{(nmesh\ ×\ G)} matrix of sample sizes
  generated by \texttt{draw\_nsamp()} using:

  \begin{itemize}
  \tightlist
  \item
    \texttt{dist\_code}, \texttt{mean\_nsamp}
  \item
    plus \texttt{ln\_sd} or \texttt{nb\_size} or \texttt{Nmin/Nmax}
    depending on \texttt{dist\_code}.
  \end{itemize}
\end{itemize}

\subsubsection{\texorpdfstring{3.3 Simulated composition counts
(\texttt{X})}{3.3 Simulated composition counts (X)}}\label{simulated-composition-counts-x}

\begin{itemize}
\tightlist
\item
  \textbf{Phi calculation}: \texttt{phi} is calculated as the inverse of
  the overdispersion parameterization used in the likelihood (e.g., for
  Dirichlet-multinomial parameterizations, typically
  \texttt{phi\ =\ 1\ /\ theta} or via model-specific scaling inside the
  MLE block).
\end{itemize}

For each simplex point \texttt{p\_true} and each group \texttt{g=1..G},
the script simulates:

\begin{enumerate}
\def\labelenumi{\arabic{enumi}.}
\tightlist
\item
  Group concentration (Dirichlet total mass):\\
  \textbf{α₀,g = (theta\_true · od\_mult) · N\_g}
\item
  Perturbed group composition on the simplex:\\
  \textbf{log(p\_g) = log(p\_true) + Normal(0, sigma)}
\item
  Dirichlet draw for multinomial probability \texttt{phi}, then
  multinomial counts:\\
  \textbf{X{[}g,{]} \textasciitilde{} Multinomial(N\_g, phi)}
\end{enumerate}

These simulated counts \texttt{X} are the data used by all MLE blocks
(i)--(v).

\begin{center}\rule{0.5\linewidth}{0.5pt}\end{center}

\subsection{4) Estimated-parameter mapping for models
(i)--(v)}\label{estimated-parameter-mapping-for-models-iv}

\subsubsection{4.1 MLE Calculation Line Numbers by
Model}\label{mle-calculation-line-numbers-by-model}

\begin{itemize}
\tightlist
\item
  Line 440:
  \texttt{fit\_i\ \textless{}-\ optim(t0\_i,\ nll\_i,\ score\_param\_i,\ method\ =\ "BFGS",}
\item
  Line 464:
  \texttt{fit\_ii\ \textless{}-\ try(optim(c(z0,\ t0),\ nll\_ii,\ grad\_ii,\ method\ =\ "BFGS",}
\item
  Line 486:
  \texttt{fit\_iii\ \textless{}-\ try(optim(z0,\ nll\_iii,\ grad\_iii,\ method\ =\ "BFGS",}
\item
  Line 507:
  \texttt{fit\_v\ \textless{}-\ try(optim(c(z0,\ 0),\ nll\_v,\ grad\_v,\ method\ =\ "BFGS",}
\end{itemize}

All models produce an estimate of the \textbf{shared composition vector}
\texttt{p} (length \texttt{K}) using a softmax parameterization:

\begin{itemize}
\tightlist
\item
  \textbf{Unconstrained} parameter: \texttt{z} (length \texttt{K-1})
\item
  \textbf{Simplex} parameter: \texttt{p\ =\ softmax\_from\_z(z)}
\end{itemize}

\subsubsection{\texorpdfstring{Model (i): shared unconstrained
\texttt{alpha\_vec}}{Model (i): shared unconstrained alpha\_vec}}\label{model-i-shared-unconstrained-alpha_vec}

\textbf{Estimated parameters} - \texttt{alpha\_vec} (length \texttt{K},
positive), shared across all groups

\textbf{Optimization parameterization} - Optimized parameter:
\texttt{t\ =\ log(alpha\_vec)} (length \texttt{K}) - After optimization:
\texttt{alpha\_hat\_i\ =\ exp(fit\_i\$par)} - Implied simplex estimate:
\texttt{p\_hat\_i\ =\ alpha\_hat\_i\ /\ sum(alpha\_hat\_i)}

\textbf{Where in script} - Block header: \texttt{\#\#\ MLE\ (i)} (around
line \textbf{421})

\begin{center}\rule{0.5\linewidth}{0.5pt}\end{center}

\subsubsection{\texorpdfstring{Model (ii): shared simplex \texttt{p} and
group-specific total mass
\texttt{a0\_g}}{Model (ii): shared simplex p and group-specific total mass a0\_g}}\label{model-ii-shared-simplex-p-and-group-specific-total-mass-a0_g}

\textbf{Estimated parameters} - \texttt{p} (length \texttt{K}, simplex)
- \texttt{a0\_g} for each group \texttt{g} (length \texttt{G}, positive)

\textbf{Optimization parameterization} - Optimized parameter vector:
\texttt{par\ =\ c(z,\ t)} - \texttt{z} length \texttt{K-1} →
\texttt{p\ =\ softmax\_from\_z(z)} - \texttt{t} length \texttt{G} →
\texttt{a0\ =\ exp(t)}

\textbf{Where in script} - Block header: \texttt{\#\#\ MLE\ (ii)}
(around line \textbf{447})

\begin{center}\rule{0.5\linewidth}{0.5pt}\end{center}

\subsubsection{\texorpdfstring{Model (iii): shared simplex \texttt{p},
with fixed linear scaling α₀,g = θ ·
N\_g}{Model (iii): shared simplex p, with fixed linear scaling α₀,g = θ · N\_g}}\label{model-iii-shared-simplex-p-with-fixed-linear-scaling-ux3b1ux2080g-ux3b8-n_g}

\textbf{Estimated parameters} - \texttt{p} (length \texttt{K}, simplex)

\textbf{Fixed inputs used} -
\texttt{theta\_fixed\ \textless{}-\ theta\_true} (θ is fixed, not
estimated) - \texttt{N\_vec{[}g{]}} (sample sizes)

\textbf{Model structure} - α₀,g = θ\_fixed · N\_g

\textbf{Optimization parameterization} - Optimize \texttt{z} (length
\texttt{K-1}), map to \texttt{p\ =\ softmax\_from\_z(z)}

\textbf{Where in script} - Block header: \texttt{\#\#\ MLE\ (iii)}
(around line \textbf{474})

\begin{center}\rule{0.5\linewidth}{0.5pt}\end{center}

\subsubsection{\texorpdfstring{Model (iv): multinomial MLE with shared
simplex
\texttt{p}}{Model (iv): multinomial MLE with shared simplex p}}\label{model-iv-multinomial-mle-with-shared-simplex-p}

\textbf{Estimated parameters} - \texttt{p} (length \texttt{K}, simplex)

\textbf{Model structure} - Multinomial log-likelihood using pooled
counts across groups (no Dirichlet layer)

\textbf{Where in script} - Block header: \texttt{\#\#\ MLE\ (iv)}
(around line \textbf{516})

\begin{center}\rule{0.5\linewidth}{0.5pt}\end{center}

\subsubsection{\texorpdfstring{Model (v): shared simplex \texttt{p} and
estimated linear scaling θ, α₀,g = θ ·
N\_g}{Model (v): shared simplex p and estimated linear scaling θ, α₀,g = θ · N\_g}}\label{model-v-shared-simplex-p-and-estimated-linear-scaling-ux3b8-ux3b1ux2080g-ux3b8-n_g}

\textbf{Estimated parameters} - \texttt{p} (length \texttt{K}, simplex)
- \texttt{theta} (scalar positive)

\textbf{Optimization parameterization} - Optimized parameter vector:
\texttt{par\ =\ c(z,\ t)} - \texttt{z} length \texttt{K-1} →
\texttt{p\ =\ softmax\_from\_z(z)} - \texttt{t} scalar →
\texttt{theta\ =\ exp(t)}

\textbf{Where in script} - Block header: \texttt{\#\#\ MLE\ (v)} (around
line \textbf{494})

\begin{center}\rule{0.5\linewidth}{0.5pt}\end{center}

\subsection{\texorpdfstring{4) Output CSV structure (created by
\texttt{hake.R})}{4) Output CSV structure (created by hake.R)}}\label{output-csv-structure-created-by-hake.r}

When \texttt{hake.R} finishes evaluating the simplex grid, it writes a
CSV file named:

\begin{itemize}
\tightlist
\item
  \texttt{csv\_path\ \textless{}-\ paste0(out\_prefix,\ ".csv")} (around
  line \textbf{585} in \texttt{hake.R})
\end{itemize}

\subsubsection{4.1 Row definition}\label{row-definition}

Each row corresponds to \textbf{one simplex point} \texttt{p\_true} from
\texttt{mesh} (length \texttt{K}) and the associated
simulation/estimation results for that point. The script then appends
the group sample sizes \texttt{N1..NG} used for that row.

If a simplex point violates the component-wise bounds
(\texttt{p\_lbound} / \texttt{p\_ubound}), the row is still returned but
\textbf{all model outputs are set to \texttt{NA}} to avoid dropping rows
(see the early return around line \textbf{524}).

\subsubsection{4.2 Column order}\label{column-order}

For a given \texttt{K} and \texttt{G}, the CSV columns are written in
this order (see column naming around lines \textbf{552--563}, and
binding of \texttt{N1..NG} around lines \textbf{565--580}):

\begin{enumerate}
\def\labelenumi{\arabic{enumi}.}
\tightlist
\item
  \textbf{True simplex coordinates}

  \begin{itemize}
  \tightlist
  \item
    \texttt{p1,\ p2,\ ...,\ pK}
  \end{itemize}
\item
  \textbf{Model (i)}

  \begin{itemize}
  \tightlist
  \item
    \texttt{rmse\_i}\strut \\
  \item
    \texttt{p\_hat\_i\_1,\ p\_hat\_i\_2,\ ...,\ p\_hat\_i\_K}
  \end{itemize}
\item
  \textbf{Model (ii)}

  \begin{itemize}
  \tightlist
  \item
    \texttt{rmse\_ii}\strut \\
  \item
    \texttt{p\_hat\_ii\_1,\ ...,\ p\_hat\_ii\_K}
  \end{itemize}
\item
  \textbf{Model (iii)}

  \begin{itemize}
  \tightlist
  \item
    \texttt{rmse\_iii}\strut \\
  \item
    \texttt{p\_hat\_iii\_1,\ ...,\ p\_hat\_iii\_K}
  \end{itemize}
\item
  \textbf{Model (iv)}

  \begin{itemize}
  \tightlist
  \item
    \texttt{rmse\_iv}\strut \\
  \item
    \texttt{p\_hat\_iv\_1,\ ...,\ p\_hat\_iv\_K}
  \end{itemize}
\item
  \textbf{Model (v)}

  \begin{itemize}
  \tightlist
  \item
    \texttt{rmse\_v}\strut \\
  \item
    \texttt{p\_hat\_v\_1,\ ...,\ p\_hat\_v\_K}
  \end{itemize}
\item
  \textbf{Group sample sizes appended at the end}

  \begin{itemize}
  \tightlist
  \item
    \texttt{N1,\ N2,\ ...,\ NG}
  \end{itemize}
\end{enumerate}

\subsubsection{4.3 Interpretation notes}\label{interpretation-notes}

\begin{itemize}
\tightlist
\item
  Each \texttt{p\_hat\_*} is an estimated simplex vector (length
  \texttt{K}) produced by the corresponding MLE block.
\item
  \texttt{rmse\_*} is the RMSE between \texttt{p\_hat\_*} and
  \texttt{p\_true} for that row.
\item
  The CSV contains \textbf{both} the truth (\texttt{p1..pK}) and each
  model's paired \texttt{(rmse,\ p\_hat)} output, so downstream analysis
  can avoid re-running the MLEs.
\end{itemize}

\subsection{5) Practical notes for keeping the mapping
``current''}\label{practical-notes-for-keeping-the-mapping-current}

If you add new keys to \texttt{hake.inp}, the mapping should be updated
in \textbf{two places}:

\begin{enumerate}
\def\labelenumi{\arabic{enumi}.}
\tightlist
\item
  The typed assignment block where \texttt{kv{[}{[}"..."{]}{]}} gets
  converted (near lines \textasciitilde244--290).
\item
  The \texttt{params\ \textless{}-\ list(...)} block (near lines
  \textasciitilde300--315), so the run prints a complete record of
  inputs.
\end{enumerate}

\begin{center}\rule{0.5\linewidth}{0.5pt}\end{center}

\subsection{\texorpdfstring{Appendix A: Example \texttt{hake.inp} keys
(for
reference)}{Appendix A: Example hake.inp keys (for reference)}}\label{appendix-a-example-hake.inp-keys-for-reference}

This is the \textbf{expected} set of keys (order does not matter), using
key=value format:

\begin{itemize}
\tightlist
\item
  \texttt{K}, \texttt{G}, \texttt{h}
\item
  \texttt{theta\_true}, \texttt{od\_mult}, \texttt{sigma}
\item
  \texttt{nsims}, \texttt{random.seed}
\item
  \texttt{dist\_code}, \texttt{mean\_nsamp}, plus \texttt{ln\_sd} or
  \texttt{nb\_size} or \texttt{Nmin}/\texttt{Nmax} depending on
  \texttt{dist\_code}
\item
  Optional: \texttt{p\_lbound}, \texttt{p\_ubound}
\end{itemize}

\end{document}
